\section{Local cube laws and lifts of face symmetries}
\label{sec:cube-limits}

The comparison from Section~\ref{sec:rough-progressions} will give symmetries
of individual marginals of a joint nilmanifold-valued state.  We need to act on
the joint state even when those symmetries do not preserve its other
marginals.  The purpose of this Section is to construct such actions on a
universal cover.  Their nilpotence class, rather than the dimension of the
cover, is what will control the finite recurrence argument.
For the cube formalism in higher-order ergodic theory, see \cite{HostKra}.
We prove the algebraic and lifting statements in the precise forms used here.

Throughout this Section, a rational subgroup of a connected simply connected
nilpotent Lie group means a connected subgroup whose Lie algebra is rational
for the rational structure determined by the specified lattice.  We use the
basic rational Lie theory of nilmanifolds: such a subgroup is closed, its
lattice intersection is cocompact, and its orbit in the nilmanifold is an
embedded compact homogeneous space.  A rational element is an element with
rational logarithm in this structure.  Conjugation by a rational element
preserves the rational structure.  For these rational-coordinate facts
and the periodicity of rational polynomial sequences, see
\cite[Section~2 and Appendix~A]{GreenTao}.
All filtrations below have connected
rational terms and satisfy
\begin{equation}
 H_0=H_1=H\supseteq H_2\supseteq\cdots\supseteq H_s
 \supseteq H_{s+1}=\{1\},\qquad
 [H_i,H_j]\subseteq H_{i+j}.
 \label{eq:cube-filtration}
\end{equation}
Terms beyond $s$ are trivial.  We write $\mathfrak h_i=\log H_i$.
In particular $H$ has nilpotence class at most $s$.

\subsection{Upper-face groups and the missing corner}

For $q\geq0$, identify the vertices of a $q$-cube with subsets of $[q]$.
For $D\subseteq[q]$, let $g_D$ be the group-valued array equal to $g$ at
vertices containing $D$, and equal to the identity elsewhere.  The case
$D=\varnothing$ is the constant array.  Define
\[
 C^q(H_\bullet)=\HK^q(H_\bullet)
 =\left\langle g_D:\ D\subseteq[q],\quad g\in H_{|D|}\right\rangle
 \subseteq H^{\{0,1\}^q}.
\]
For $q=0$ this is $H$.

\begin{lemma}[Upper-face coordinates]\label{lem:upper-face-coordinates}
The group $C^q(H_\bullet)$ is connected, simply connected, and rational.
Order the subsets of $[q]$ by increasing cardinality, breaking ties in any
fixed way.  Every element of $C^q(H_\bullet)$ has a unique expression
\begin{equation}
 \prod_{D\subseteq[q]}(g_D)_D,\qquad g_D\in H_{|D|},
 \label{eq:ordered-face-product}
\end{equation}
in this order; the map from the face coefficients to the array is a
diffeomorphism.  An array in this group which is the identity except possibly
at the all-ones vertex has its remaining entry in $H_q$.
The group is invariant under permutations and reflections of cube
coordinates.  If $\Gamma$ is a lattice in $H$, its image in
$(H/\Gamma)^{\{0,1\}^q}$ is compact.
\end{lemma}

\begin{proof}
For $X\in\mathfrak h_{|D|}$ write $X_D$ for the corresponding face array.
The Boolean monomials $\prod_{j\in D}e_j$ are linearly independent, so
\[
 \mathfrak c^q
 =\bigoplus_{D\subseteq[q]}
   \mathfrak h_{|D|}\prod_{j\in D}e_j
\]
is a direct sum of vector spaces.  The entrywise Lie bracket satisfies
$[X_D,Y_E]=[X,Y]_{D\cup E}$.  The filtration places this bracket in
the summand indexed by $D\cup E$.  This remains true when a set is empty,
because $H_0=H_1$ and the filtration is decreasing.  Thus
$\mathfrak c^q$ is a rational Lie subalgebra.

Every suffix in the chosen order of summands is an ideal: a bracket with a
summand indexed by $E$ stays at $E$ or moves to a strictly larger set.
Successively passing to the quotients by these suffix ideals gives
\eqref{eq:ordered-face-product}.  Equivalently, the coefficient at $E$ is
recovered from the entry at vertex $E$ after the coefficients at proper
subsets of $E$ have been removed.  This also proves uniqueness and smoothness
of the inverse.  The connected subgroup with Lie algebra $\mathfrak c^q$
is closed and simply connected, as follows from the exponential
diffeomorphism for a simply connected nilpotent group.  It is exactly the
group generated by the faces.

The same recovery at successive vertices proves the assertion about an array
supported at the top vertex.  Reflecting a coordinate replaces $e_j$ by
$1-e_j$.  The resulting face indicator is a linear combination of upper-face
indicators of no greater cardinality; the decreasing filtration therefore
preserves $\mathfrak c^q$, and hence its exponential group.  Permutations are
immediate.  Finally rationality gives a cocompact lattice
$C^q(H_\bullet)\cap\Gamma^{\{0,1\}^q}$, so the indicated image is compact.
\end{proof}

We will also use the following explicit rational filtration on the cube
group:
\begin{equation}
 \mathfrak c^q_i
 =\bigoplus_{D\subseteq[q]}
   \mathfrak h_{\max(i,|D|)}\prod_{j\in D}e_j,
 \qquad i\geq0.
 \label{eq:cube-group-filtration}
\end{equation}
It has $C^q_0=C^q_1=C^q(H_\bullet)$ and degree at most $s$.
Indeed the bracket of its $D$ and $E$ summands lies at filtration level at
least both $i+j$ and $|D\cup E|$.  Thus it satisfies the required bracket
inclusions.

The following missing-corner constraint and its coordinate-product
consequence appear in \cite[Proposition~11.5, proof of
Proposition~11.2, and Appendix~E]{GreenTaoLinear}.  We give the
face-coordinate proof, including the uniform approximation over
bounded observable families needed in Section~\ref{sec:prediction}.

\begin{lemma}[Continuous missing-corner reconstruction]
\label{lem:cube-corner}
Let $G$ be a connected simply connected nilpotent Lie group of step at most
$s$, and let $\Gamma$ be a lattice in $G$. Equip $G$ with its lower central
filtration, with $G_0=G_1=G$.
In dimension $s+1$, any $2^{s+1}-1$ vertices of a cube in the image of
$C^{s+1}(G_\bullet)$ determine the last vertex uniquely and continuously
on the set of admissible partial cubes.  Linear orbit cubes
\[
 \bigl(g^{b_0+\sum_{j=1}^{s+1}e_jb_j}x\bigr)_e
\]
belong to this compact cube set.

Consequently, for every $\eps>0$ and every bounded Lipschitz observable $F$
on $G/\Gamma$, its value at the missing vertex is uniformly within $\eps$
of a finite sum of products of bounded continuous functions of the other
individual vertices.  Those individual functions may be taken Lipschitz
and bounded by one, with scalar coefficients outside the products.  For a
fixed finite list of nilmanifolds and fixed bounds on $\norm{F}_\infty$ and
$\Lip(F)$, finitely many such approximation recipes suffice at each
accuracy.
\end{lemma}

\begin{proof}
Reflect the missing vertex to the all-ones vertex.  If two group cubes
$a,b\in C^{s+1}(G_\bullet)$ agree modulo $\Gamma$ at all other vertices,
then $z=a^{-1}b$ is a group cube whose entries at all proper subsets of
$[s+1]$ belong to $\Gamma$.  Recover its face coefficients by
\eqref{eq:ordered-face-product}.  Inductively the coefficient at a proper
subset $E$ lies in $\Gamma$, because its value is the entry at $E$
multiplied by inverses of previously recovered lattice elements.  The top
coefficient belongs to $G_{s+1}=\{1\}$.  The top entry of $z$ is therefore
also a product of lattice elements.  The original two top vertices agree
modulo $\Gamma$.

The deletion map from the compact cube set to the space of partial cubes
is now a continuous injection into a Hausdorff space.  It is a homeomorphism
onto its image, proving continuity of reconstruction.  For a representative
$x_0\in G$ of $x$, the displayed linear cube is represented by the product
of the constant arrays $g^{b_0}$ and $x_0$, with the codimension-one
arrays $g^{b_j}$ between them.  It therefore lies in the cube group.

On the compact image of the deletion map, products of continuous functions
of individual coordinates form an algebra containing constants and
separating points.  Stone--Weierstrass approximates the reconstructed
observable by finite sums from this algebra.  Lipschitz functions are dense
in the continuous functions on each compact nilmanifold, so the factors
can be made Lipschitz and then normalized to have sup norm at most one.
Finally the family of observables with the prescribed sup and Lipschitz
bounds is compact in the uniform norm.  A finite uniform net, followed by
the construction for each member of that net, proves the last assertion.
\end{proof}

\subsection{Adapted factorization on a sequence of scales}

The lifting argument will use face information in every dimension
through $s+1$.  We therefore need one rational filtration whose cube
groups describe all these local laws, together with the point law.
To obtain it, we first factor the orbit.  The left factor will vary
on the scale of the full interval, the right factor will be periodic
modulo the lattice, and the remaining factor will lie in a fixed
rational subgroup.  Its Taylor coefficients will satisfy the
irrationality conditions that give equidistribution in all the
associated cube groups.
The smooth--equidistributed--periodic decomposition of Green and Tao
\cite[Theorem~1.19]{GreenTao} was refined using Taylor-coefficient
irrationality and filtration descent in
\cite[Definitions~A.5--A.6 and Lemma~2.9]{GreenTaoRegularity}.
We prove the sequential limiting form needed for the joint cube laws.

A polynomial $P:\R\to H$ is \emph{adapted} to $H_\bullet$ if
\[
 \log P(b)=\sum_{j=0}^s b^j U_j,\qquad U_j\in\mathfrak h_j.
\]
This definition is unchanged on passing from monomials to binomial
polynomials.  Products and inverses of adapted polynomials are adapted:
in the finite Baker--Campbell--Hausdorff expansion, a term of polynomial
degree $j$ belongs to $\mathfrak h_j$.  Terms of degree above $s$ vanish.
Constant factors cause no difficulty, since $H_0$ normalizes every level.
This definition implies polynomiality by group differences, as required
in Theorem~\ref{thm:quantitative-leibman}.  More explicitly, suppose the
coefficient of degree $j$ in $\log R(b)$ lies in
$\mathfrak h_{j+v}$.  Give a formal shift $t$ degree one as well.
BCH shows that every coefficient of total degree $j+k$ in
$\log(R(b+t)R(b)^{-1})$ lies in $\mathfrak h_{j+k+v}$; brackets of
two terms can only increase this filtration index.  The expression is
zero at $t=0$, so every surviving monomial uses at least one power of
$t$.  For each fixed $t$, its degree-$j$ coefficient in $b$ therefore
belongs to $\mathfrak h_{j+v+1}$.  Iterating from $v=0$ puts every
$k$-fold difference in $H_k$.  The same argument works with total
degree in several variables.

\begin{lemma}[Taylor coordinates and finite descent]
\label{lem:adapted-factorization}
Fix a connected simply connected nilpotent group $K$, a lattice $\Lambda$,
and a rational filtration $K_\bullet$ satisfying
\eqref{eq:cube-filtration}.  Let $L\to\infty$ along any sequence, and let
$P_L$ be arbitrary adapted polynomials.  After passing to a subsequence
there exist a fixed rational filtration $H_\bullet$ with $H_i\subseteq K_i$,
a positive integer $d_0$, rational elements $\sigma_0,\ldots,\sigma_{d_0-1}$
of $K$, and factorizations
\begin{equation}
 P_L(b)\Lambda
 =h_L(b)Q_L(b)\gamma_L(b)\Lambda,
 \qquad
 Q_L(b)=\prod_{j=1}^s a_{j,L}^{\binom bj},\quad a_{j,L}\in H_j,
 \label{eq:adapted-factorization}
\end{equation}
such that the following hold.
\begin{enumerate}
\item The maps $\beta\mapsto h_L(L\beta)$ converge locally uniformly on
$\R$ to a smooth map $h:\R\to K$.  They and their first derivatives are
bounded on every fixed compact interval, uniformly in $L$.
\item For integer $b$,
$\gamma_L(b)\Lambda=\sigma_{b\bmod d_0}\Lambda$.
\item If $1\leq j\leq s$ and $\xi:H_j\to\R$ is a nonzero rational
homomorphism annihilating $H_{j+1}$ and every $[H_i,H_{j-i}]$,
$1\leq i<j$, then
\begin{equation}
 L^j\norm{\xi(a_{j,L})}_{\R/\Z}\longrightarrow\infty.
 \label{eq:level-irrationality}
\end{equation}
Here rational homomorphisms include all nonzero rational multiples of one
another.  No boundedness of the Taylor coefficients $a_{j,L}$ is asserted
or required.
\end{enumerate}
\end{lemma}

\begin{proof}
We first justify Taylor coordinates.  For any adapted $R$ write
$a_0=R(0)$ and remove this constant on the left.  Suppose that, after
removing the factors of orders below $j$, the remaining adapted polynomial
$R_j$ is the identity at $0,\ldots,j-1$.  The $j$th forward difference of
$\log R_j$ at zero equals $\log R_j(j)$: all its other evaluated terms
vanish.  Terms of degree less than $j$ in the logarithm make no
contribution, and all terms of degree at least $j$ have coefficients in
$\mathfrak h_j$.  Thus $a_j=R_j(j)\in H_j$.  Multiplication on the left by
$a_j^{-\binom bj}$ leaves an adapted polynomial vanishing at the first
$j+1$ integers.  After order $s$, its logarithm has degree at most $s$
and $s+1$ zeros, so it is identically zero.  This proves the ordered Taylor
representation.

Choose $\lambda_L\in\Lambda$ so that
$c_L=P_L(0)\lambda_L$ lies in a fixed compact set of representatives.
Normalize by
$R_L(b)=c_L^{-1}P_L(b)\lambda_L$.
It is adapted, has constant term the identity, and
$P_L(b)\Lambda=c_LR_L(b)\Lambda$.  Adaptation is preserved here either
by the preceding BCH observation, or by noting that conjugation by
$\lambda_L$ preserves all terms of $K_\bullet$.

We describe one descent step for a fixed current rational filtration
$H_\bullet$ and a normalized adapted polynomial with Taylor coefficients
$a_{j,L}$.  If \eqref{eq:level-irrationality} fails, fix one of its
characters $\xi$ and pass to a subsequence on which
\[
 \xi(a_{j,L})=m_L+\lambda'_L,\qquad
 m_L\in\Z,\qquad \abs{\lambda'_L}\leq C L^{-j}.
\]
Choose a fixed rational $V\in\mathfrak h_j$ with $\xi(\exp V)=1$.
Replace the polynomial $R_L$ by
\begin{equation}
 R'_L(b)
 =\exp\bigl(-\lambda'_L\tbinom bj V\bigr)
   R_L(b)
   \exp\bigl(-m_L\tbinom bj V\bigr).
 \label{eq:filtration-descent}
\end{equation}
It remains adapted to the old filtration.  Its Taylor coefficients of
orders less than $j$ are unchanged, since both multipliers are the
identity at $0,\ldots,j-1$.  Modulo $H_{j+1}$, the order-$j$ coefficient
is the old coefficient with the two displayed corrections.  To verify this
even in the presence of lower Taylor factors, commute those factors past
the left correction at $b=j$; the commutators lie in
$[H,H_j]\subseteq H_{j+1}$.  Its $\xi$-value is consequently
$-\lambda'_L+\xi(a_{j,L})-m_L=0$.

Replace $H_j$ by $H_j\cap\ker\xi$ and leave all other positive levels
unchanged.  When $j=1$ replace $H_0$ by the same kernel.  The kernel is a
connected rational subgroup, since in exponential coordinates it is the
kernel of a rational linear form.  Nesting is preserved because $\xi$
kills $H_{j+1}$.  A bracket inclusion whose target is the changed level
is preserved because $\xi$ kills each $[H_i,H_{j-i}]$.  Inclusions whose
target has larger index already land in $H_{j+1}$, and the other
inclusions are unchanged or have a smaller source.  Thus this is again a
filtration.  The Taylor calculation shows that $R'_L$ is adapted to it.
In the case $j=1$, all higher coefficients already belong to $H_2$,
so its entire image lies in the new $H$.

The removed factors give
\[
 R_L(b)=
 \exp\bigl(\lambda'_L\tbinom bj V\bigr)
 R'_L(b)
 \exp\bigl(m_L\tbinom bj V\bigr).
\]
The left factor is bounded and has derivative $O(L^{-1})$ on every
interval $\abs b\leq A L$, with constants allowed to depend on $A,C,V$.
The right factor is an integer power of the fixed rational element
$\exp V$ at integer inputs.  Repeating the procedure strictly reduces
$\sum_{i=1}^s\dim H_i$ at every step.  It therefore terminates after
finitely many steps.  If any fixed rational level character still failed
\eqref{eq:level-irrationality} on the final subsequence, the same operation
would give another strict reduction.  Hence all the asserted character
limits hold.

There is no conjugation of accumulated smooth factors by rational factors
in this procedure: from
$P=\epsilon_1R_1\gamma_1$ and
$R_1=\epsilon_2R_2\gamma_2$ one gets
$P=(\epsilon_1\epsilon_2)R_2(\gamma_2\gamma_1)$.
Thus $h_L$ is the product of $c_L$ and finitely many of the bounded smooth
factors just described.  Pass to a subsequence on which $c_L$ converges
and every bounded scalar $L^j\lambda'_L$ converges.  Since
$L^{-j}\binom{L\beta}{j}$ converges with all derivatives on compact sets
to $\beta^j/j!$, this gives the first conclusion.

For completeness, the right factors have a common bounded period modulo
$\Lambda$, even though the integers $m_L$ need not be bounded.  Finitely
many fixed rational elements, together with a finite generating set of
$\Lambda$, generate a discrete group $\Lambda'$ containing $\Lambda$.
Here is the denominator argument.  Collect a word into integer powers of
the finitely many generator and iterated-commutator types of length at most
the nilpotence class.  Collection terminates because a commutation error
has strictly greater commutator length.  The logarithms of these finitely
many types have rational coordinates.  Applying the finite BCH formula to
their ordered powers bounds every coordinate denominator by one fixed
integer, independent of the word and its integer exponents.  This proves
discreteness.  A compact fundamental set for $\Lambda$ then shows that
$[\Lambda':\Lambda]<\infty$.

The normal core $\Lambda''=\bigcap_{g\in\Lambda'}g\Lambda g^{-1}$ has
finite index in $\Lambda'$, so $\Lambda'/\Lambda''$ is a finite group.
Let $a$ be a common multiple of its element orders.  Each
$\binom bj$ modulo $a$ has a fixed period, for example $a s!$ for
$0\leq j\leq s$: Vandermonde's identity shows that
$\binom{b+a s!}{j}-\binom bj$ is divisible by $a$.
Consequently every accumulated right factor is periodic in this finite
quotient with one common period $d_0$.  There are only finitely many
possible maps from its residue classes to $\Lambda'/\Lambda''$.
Pass to a subsequence on which this map is fixed and choose rational
representatives $\sigma_r$.  This gives the second conclusion and
completes the proof.
\end{proof}

\subsection{The joint local point and cube laws}

The factorization has fixed a subgroup, a filtration, and a finite
period.  We now identify the distribution of the remaining factor
and its cubes.  The key step is to detect any horizontal obstruction
in a cube group in one of the Taylor coefficients controlled above.

For a rational element $\sigma\in K$, put
\[
 \Gamma_\sigma=H\cap\sigma\Lambda\sigma^{-1},\qquad
 \Gamma^{[q]}_\sigma
 =C^q(H_\bullet)\cap
     (\sigma\Lambda\sigma^{-1})^{\{0,1\}^q}.
\]
These are lattices in their respective groups.  By Haar measure on
$hH\sigma\Lambda/\Lambda$ we mean the image of invariant probability on
$H/\Gamma_\sigma$ under $y\mapsto hy\sigma\Lambda$.  Use the analogous
convention for cube cosets.

\begin{proposition}[Simultaneous local cube law]
\label{prop:local-cube-law}
Under the hypotheses of Lemma~\ref{lem:adapted-factorization}, retain the
subsequence and its data.  Define probability measures
\begin{equation}
 m_{\beta,r}^{[q]}
 =\text{Haar image on }
 h(\beta)^{\mathrm{diag}} C^q(H_\bullet)
 \sigma_r^{\mathrm{diag}}\Lambda^{\{0,1\}^q}/
              \Lambda^{\{0,1\}^q}.
 \label{eq:local-cube-haar}
\end{equation}
The empirical point laws satisfy
\begin{equation}
 \frac1{\lfloor L\rfloor}\sum_{0\leq b<\lfloor L\rfloor}
       \delta_{P_L(b)\Lambda}
 \ \Longrightarrow
 \frac1{d_0}\sum_{r=0}^{d_0-1}\int_0^1 m_{\beta,r}^{[0]}\,d\beta.
 \label{eq:joint-point-law}
\end{equation}
For every fixed cube dimension $q$, every $\beta\in(0,1)$, every residue
$r\bmod d_0$, and every bounded Lipschitz test $F$ on
$(K/\Lambda)^{\{0,1\}^q}$,
\begin{equation}
 \lim_{\alpha\downarrow0}\limsup_{L\to\infty}
 \left|
 \E_{\substack{b_0/L\in[\beta,\beta+\alpha),\ b_0\equiv r\ (d_0)\\
                 b_j/L\in[0,\alpha),\ b_j\equiv0\ (d_0),\ 1\leq j\leq q}}
 F\bigl((P_L(b_0+\sum_{j=1}^q e_jb_j)\Lambda)_e\bigr)
 -\int F\,dm_{\beta,r}^{[q]}
 \right|=0.
 \label{eq:local-cube-limit-order}
\end{equation}
The inputs in this average are independent and uniform on the specified
integer progressions.  The same assertion holds for boxes with side
lengths between $c\alpha L$ and $C\alpha L$, whose base coordinates
are within $C\alpha L$ of $\beta L$ and whose increments have
absolute value at most $C\alpha L$, where $c,C>0$ are fixed before
either limit.
For fixed bounds the assertions are uniform over the tests and these
boxes.  In particular, arbitrary fixed-modulus progression restrictions
are permitted if the base has residue $r$ and all increments have residue
zero modulo $d_0$.  Every progression modulus is fixed before the
length limit, which always precedes the locality limit.
The same factorization supplies these conclusions for every fixed $q$.
\end{proposition}

\begin{proof}
We first prove equidistribution before inserting the smooth factor.  Fix
$q,r$ and abbreviate $C=C^q(H_\bullet)$.  The cube polynomial
\[
 \mathbf Q_L(\mathbf b)
 =\bigl(Q_L(b_0+\textstyle\sum_{j=1}^q e_jb_j)\bigr)_e
\]
takes values in $C$.  To see this and record the needed coefficient
information, expand
\begin{equation}
 \binom{b_0+\sum_{j=1}^q e_jb_j}{\ell}
 =\sum_{D\subseteq[q]}
      \Bigl(\prod_{j\in D}e_j\Bigr)F_{\ell,D}(\mathbf b).
 \label{eq:boolean-binomial-expansion}
\end{equation}
Each $F_{\ell,D}$ has degree at most $\ell$, and every one of its
monomials uses exactly the increment variables indexed by $D$, each with
positive exponent.  In particular it vanishes for $|D|>\ell$.
Thus the $\ell$th Taylor factor on the cube is the product of the face
arrays $(a_{\ell,L})_D^{F_{\ell,D}(\mathbf b)}$ with $|D|\leq\ell$.
These are in $C$, because $H_\ell\subseteq H_{|D|}$.  This expansion
also shows adaptation to \eqref{eq:cube-group-filtration}: a coefficient
of total polynomial degree $k$ in a factor of order $\ell\geq k$ lies
in $\mathfrak h_\ell\subseteq\mathfrak h_{\max(k,|D|)}$.
BCH preserves this assertion on multiplying the Taylor factors.

We claim that $\mathbf Q_L$ is equidistributed on
$C/\Gamma^{[q]}_{\sigma_r}$ on every box whose sides are fixed positive
fractions of $L$, whose origin is $O(L)$, and with any fixed progression
restrictions.  The error tends to zero uniformly for fixed bounds on
these data and on the Lipschitz tests.  Otherwise
Theorem~\ref{thm:quantitative-leibman}, applied to the fixed rational
nilmanifold and the filtration \eqref{eq:cube-group-filtration}, supplies
along a subsequence a nonzero horizontal character from a fixed finite
list.  Pass to a further subsequence on which this character, denoted
$\chi:C\to\R$, is fixed.  It annihilates commutators and takes integer
values on $\Gamma^{[q]}_{\sigma_r}$.

For a face $D$, write $\chi_D(y)=\chi(y_D)$ wherever this is defined.
Choose the largest $j\geq1$ such that for some $|D|\leq j$ the
restriction
$\xi=\chi_D|_{H_j}$ is nonzero.  Such a pair exists because the face
subgroups generate $C$, including the empty face and $H_0=H_1$.
The restriction is a rational homomorphism: it is integer-valued on the
lattice $H_j\cap\sigma_r\Lambda\sigma_r^{-1}$, and conjugation by
$\sigma_r$ preserves rationality.  Maximality makes it vanish on
$H_{j+1}$.  For $1\leq i<j$, choose $A,B\subseteq D$ with
$A\cup B=D$, $|A|\leq i$, and $|B|\leq j-i$; one or both sets may
be empty.  If $u\in H_i$ and $v\in H_{j-i}$, their face arrays on
$A,B$ lie in $C$ and
$[u_A,v_B]=[u,v]_D$.  Since $\chi$ kills this commutator, $\xi$ kills
$[H_i,H_{j-i}]$.  Thus $\xi$ is exactly a character covered by
\eqref{eq:level-irrationality}.

Applying $\chi$ to \eqref{eq:boolean-binomial-expansion} gives the
scalar polynomial
\[
 \chi(\mathbf Q_L(\mathbf b))
 =\sum_{\ell,D}\chi_D(a_{\ell,L})F_{\ell,D}(\mathbf b).
\]
All terms with $\ell>j$ vanish by maximality.  In its degree-$j$ part,
the coefficient of
\begin{equation}
 b_0^{j-|D|}\prod_{k\in D}b_k
 \quad\hbox{is}\quad
 \frac{\xi(a_{j,L})}{(j-|D|)!}.
 \label{eq:isolated-horizontal-coefficient}
\end{equation}
Indeed the top part of the binomial polynomial is the $j$th power
divided by $j!$.  No other face can contribute to this monomial: its
increment support must be exactly $D$.  No lower Taylor order has
degree $j$.  This proves both the value and the absence of cancellation
in \eqref{eq:isolated-horizontal-coefficient}.

On a fixed residue progression, substitute
$b_\nu=r_\nu+d t_\nu$ with fixed $d$.  Its top-degree coefficient is
multiplied by $d^j$; translations of the box origin do not alter it.
The bounded scalar coefficient obstruction in
Theorem~\ref{thm:quantitative-leibman}, clearing the fixed
binomial-to-monomial denominators if necessary, would therefore give
\[
 L^j\norm{a\,\xi(a_{j,L})}_{\R/\Z}=O(1)
\]
for a fixed nonzero rational number $a$.  This contradicts
\eqref{eq:level-irrationality} for the rational character $a\xi$.
The claim follows.  The argument also covers $q=0$, when only the
empty face occurs.  If $C$ is trivial, the claim is immediate.

In the averages in \eqref{eq:local-cube-limit-order}, all vertex inputs
have the same residue $r\bmod d_0$.  Consequently their right factors
can be replaced exactly by $\sigma_r\Lambda$.  For fixed $\alpha$,
the preceding claim gives the Haar law for the $Q_L$ cube.  At every
vertex, $b/L=\beta+O_q(\alpha)$; Lemma~\ref{lem:adapted-factorization}
therefore gives
\[
 h_L(b)=h(\beta)+O(\alpha)+o_{L\to\infty}(1)
\]
in any fixed local metric, uniformly on the sampled box.  Multiplication
by these bounded elements and passage to the compact nilmanifold changes
a bounded Lipschitz test by $O_F(\alpha)+o(1)$, uniformly over the
remaining entries.  The latter uniformity follows from smoothness of
the action on the compact nilmanifold.  First letting $L\to\infty$
and then $\alpha\downarrow0$ proves
\eqref{eq:local-cube-limit-order}, including its stated uniform versions.

Finally partition $[0,1]$ into finitely many intervals of length at most
$\alpha$.  On each interval and each residue class, apply the same
argument with $q=0$ and freeze $h$ at an endpoint.  The proportions of
the residue classes tend to $1/d_0$, and the interval proportions tend
to their lengths.  Letting $\alpha\downarrow0$ gives the Riemann integral
in \eqref{eq:joint-point-law}.  This integral is well defined because
the Haar images depend continuously on $\beta$.  No further subsequence
depending on $q$ was used: \eqref{eq:level-irrationality} proves the
equidistribution assertion for each fixed dimension.
\end{proof}

\subsection{Lifting a symmetry of one marginal}

The local cube law describes the joint state as a mixture of Haar
measures on compact orbits, with one filtration describing the cubes
on each orbit.  In Section~\ref{sec:alignment}, the progression
comparison will show that an automorphism of one component preserves
the projections of these cube cosets when applied on a face.  The
next proposition converts precisely that information into a
transformation of the joint cover.  We will then show that any family
of these lifts generates a group of nilpotence class at most $s$, as
required for recurrence.

Choose linear coordinates on $\mathfrak h$ by choosing, for every $j$,
a complement of $\mathfrak h_{j+1}$ in $\mathfrak h_j$.
Assign weight $j$ to coordinates on that complement.  The weighted degree
of a monomial is the sum of its variable weights, with multiplicity;
constants have weight zero.  A \emph{polynomial shear} will mean a map
\begin{equation}
 T(u)_{j,k}=u_{j,k}+p_{j,k}(u),\qquad
 \deg_{\mathrm{wt}}p_{j,k}<j.
 \label{eq:weighted-shear}
\end{equation}
In particular the displacement at weight $j$ uses only coordinates of
smaller weight.

\begin{proposition}[A marginal face symmetry has a joint lift]
\label{prop:face-lift}
Let $K=K'\times K''$ be a connected simply connected nilpotent Lie
group with lattice $\Lambda=\Lambda'\times\Lambda''$, and let $\pi$
denote the first group projection and its induced map of nilmanifolds.
Let $H\subseteq K$ be a connected rational subgroup with a connected
rational filtration $H_\bullet$ satisfying \eqref{eq:cube-filtration}.
Let $h\in K$, and
let $\sigma\in K$ be rational.  Let $S$ be a fixed automorphism of $K'$
preserving $\Lambda'$.  Suppose that, for every $1\leq q\leq s+1$,
applying $S$ to the entries on any upper codimension-one face preserves
setwise
the projected cube coset
\begin{equation}
 \pi\left(h^{\mathrm{diag}}C^q(H_\bullet)
                \sigma^{\mathrm{diag}}\Lambda^{\{0,1\}^q}/
                   \Lambda^{\{0,1\}^q}\right).
 \label{eq:projected-face-coset}
\end{equation}
Then there is a polynomial shear $T:\mathfrak h\to\mathfrak h$ such
that, for every $u\in\mathfrak h$,
\begin{equation}
 \pi\bigl(h\exp(Tu)\sigma\Lambda\bigr)
 =S\bigl(\pi(h\exp(u)\sigma\Lambda)\bigr).
 \label{eq:joint-lift-identity}
\end{equation}
The lift is a bijection of the cover.  Its action on the other component
is unrestricted.  This conclusion applies separately to every point
coset from Proposition~\ref{prop:local-cube-law}; no measurable choice
of lifts as a function of $\beta,r$ is required.
\end{proposition}

\begin{proof}
Put $J_i=\pi(H_i)$ and $J=J_1$.  These are connected rational groups;
their Lie algebras are rational images of those of $H_i$.  Their
filtration satisfies \eqref{eq:cube-filtration}, and projection of the
face generators shows
$\pi(C^q(H_\bullet))=C^q(J_\bullet)$.
In the remainder of the proof use $h,\sigma,\Lambda$ for the projected
elements and lattice.  The projected point coset
$hJ\sigma\Lambda/\Lambda$ is preserved by $S$: each of its points can
occur at a vertex on the specified face of a constant cube, and face
preservation, also for the inverse map, gives equality of the point
sets.  Choose $c\in J$ such that
\begin{equation}
 S(h\sigma)\Lambda=hc\sigma\Lambda.
 \label{eq:affine-basepoint}
\end{equation}
The ambient automorphism
\[
 \Psi(y)=h^{-1}S(h)S(y)S(h)^{-1}h
\]
satisfies
$S(hy\sigma)\Lambda=h\Psi(y)c\sigma\Lambda$ for $y\in J$.
We next verify on the group cover that $\Psi(J)=J$.

The map $y\mapsto y\sigma\Lambda$ on $J$ has stabilizer
$\Gamma_J=J\cap\sigma\Lambda\sigma^{-1}$.  Rationality makes
$J/\Gamma_J$ a compact embedded orbit.  Given a path $y(t)$ in $J$
from the identity to $y$, the path
$\Psi(y(t))c\sigma\Lambda$ stays in this orbit.  Lift it there starting
at $c\in J$.  It is also a lift under the ambient covering
$K'\to K'/\Lambda$ after multiplying on the right by $\sigma$.
Uniqueness of lifts under this covering forces the ambient lift
$\Psi(y(t))c$ to lie in $J$.  Thus $\Psi(J)\subseteq J$; equality
follows because its differential is injective and both connected groups
have the same dimension.  Consequently
\[
 \Phi(y)=c^{-1}\Psi(y)c
\]
is an automorphism of $J$, and the action on the point coset is represented
by the affine map $y\mapsto c\Phi(y)$.

We claim that its linear part is the identity on each associated graded
space:
\begin{equation}
 (\Phi_*-\mathrm{id})\mathfrak j_i\subseteq\mathfrak j_{i+1},
 \qquad 1\leq i\leq s.
 \label{eq:face-graded-identity}
\end{equation}
Fix $U\in\mathfrak j_i$ and work in dimension $i+1$.  Put
$\exp(tU)$ on the face with its first $i$ coordinates equal to one,
and the identity elsewhere.  This is a path $z(t)$ in
$C^{i+1}(J_\bullet)$.  Apply the affine map $c\Phi$ on the upper face
of the last coordinate.  The resulting array $z'(t)$ has its image
in the same cube coset by hypothesis.  At $t=0$ it is the constant-$c$
upper-face array, which itself belongs to $C^{i+1}(J_\bullet)$.

This quotient assertion lifts to
$z'(t)\in C^{i+1}(J_\bullet)$ for every $t$.  Indeed this cube group
is rational with lattice
$C^{i+1}(J_\bullet)\cap\Gamma_J^{\{0,1\}^{i+1}}$, so its quotient is
an embedded compact subspace of the product point orbit.  Lift the path
through that quotient beginning at $z'(0)$ and use uniqueness in the
ambient product covering, exactly as in the preceding paragraph.

Multiply $z'(t)$ on the left by the inverse constant-$c$ last-face
array, and then on the left by $z(t)^{-1}$.  The resulting group cube
is the identity except at the top vertex, whose value is
\[
 \exp(-tU)\Phi(\exp(tU)).
\]
Lemma~\ref{lem:upper-face-coordinates} puts this element in $J_{i+1}$.
Differentiating at $t=0$ proves \eqref{eq:face-graded-identity}.
In particular $\Phi$ preserves every filtration level.

Choose filtration-adapted coordinates on $\mathfrak j$.  The polynomial
\begin{equation}
 D(Y)=\log\bigl(c\Phi(\exp Y)\bigr)-Y
 \label{eq:affine-log-displacement}
\end{equation}
has weighted degree strictly less than the output weight in each
coordinate.  For its linear part this follows directly from
\eqref{eq:face-graded-identity}: an input from level $i$ moves only to
levels greater than $i$.  For the remaining terms use BCH with
$\log c$ and $\Phi_*Y$.  Every nonconstant bracket correction contains
at least one copy of the fixed vector $\log c$, which spends at least
one unit of filtration weight but contributes no polynomial degree in
$Y$.  A bracket contributing to output weight $j$ therefore has input
weighted degree at most $j-1$.  The constant term also has the required
degree.

The linear projection
$p=\pi_*:\mathfrak h\to\mathfrak j$ has a filtration-preserving linear
right inverse $R$.  To construct one, choose an adapted basis downstairs
and lift each basis vector of weight $i$ to $\mathfrak h_i$, which is
possible because $p(\mathfrak h_i)=\mathfrak j_i$.  Set
\begin{equation}
 T(u)=u+R D(pu).
 \label{eq:cover-shear-construction}
\end{equation}
This is a shear of the form \eqref{eq:weighted-shear}.  Explicitly, a
filtration-preserving linear map can send an input coordinate of weight
$i$ only to an output of weight at least $i$.  Substitution in $D$ and
then application of $R$ thus preserve the strict inequality between
input degree and output weight.  Moreover
$p(Tu)=pu+D(pu)=\log(c\Phi(\exp(pu)))$, proving
\eqref{eq:joint-lift-identity}.  Finally a map of
\eqref{eq:weighted-shear} is inverted successively in increasing
weight, so it is a polynomial bijection of the cover.
\end{proof}

\begin{remark}
Invariance of the point marginal alone would not suffice.  On
$\R^2/\Z^2$, the automorphism $S(x,y)=(x+y,y)$ preserves Haar measure
but is not a translation.  For $a=(0,1/3)$, applying $S$ to the upper
second face of the additive square $(0,a,0,a)$ gives $(0,a,0,Sa)$,
whose alternating corner sum is $Sa-a=(1/3,0)\ne0$ in the torus.
It therefore fails the two-dimensional face condition.  The analogous
test in dimension $i+1$ forces the graded identity
\eqref{eq:face-graded-identity} in the proof above.
\end{remark}

\subsection{Bounded-step shears and Haar approximation}

The lifts of different marginal symmetries may move the same coordinates.
The next Lemma controls the group they generate and gives every fixed
word in these lifts the same limiting distribution on the joint orbit.

\begin{lemma}[Nilpotence and expanding weighted boxes]
\label{lem:shear-haar}
On any finite-dimensional real vector space with coordinate weights in
$\{1,\ldots,s\}$, all the shears \eqref{eq:weighted-shear} form a group
of nilpotence class at most $s$.  This bound does not depend on the
dimension or the coefficients of the shears.  A direct product with any
group of ordinary translations has the same bound.

Let $\mathfrak h$ carry the adapted coordinates above, and let $\lambda_E$
be normalized Lebesgue measure on
\begin{equation}
 B_E=\{u:\ \abs{u_{j,k}}\leq E^j\text{ for all }j,k\},
 \qquad E\longrightarrow\infty.
 \label{eq:weighted-haar-boxes}
\end{equation}
For every fixed finite set $\mathcal T$ of shears,
\begin{equation}
 \max_{T\in\mathcal T}
   \norm{T_*\lambda_E-\lambda_E}_{\TV}\longrightarrow0.
 \label{eq:fixed-shear-tv}
\end{equation}
For every fixed $h\in K$ and rational $\sigma\in K$, the images of
$\lambda_E$ under $u\mapsto h\exp(u)\sigma\Lambda$ converge weakly to
the Haar image on $hH\sigma\Lambda/\Lambda$.  The same limit holds
after any fixed shear, and hence after each word in any fixed finite
list of words in the lifts from Proposition~\ref{prop:face-lift}.
\end{lemma}

\begin{proof}
Composition preserves \eqref{eq:weighted-shear}, and inversion does so
by solving one weight at a time.  Let $V_d$ be the space of polynomials
of weighted degree at most $d$, including the constants, and put
$V_{-1}=0$.  Substitution by a shear preserves $V_s$, and substitution
minus the identity maps $V_d$ into $V_{d-1}$.  Indeed every term in the
difference of a substituted monomial uses at least one strict
lower-degree displacement in place of its coordinate.  Use inverse
substitution if necessary to obtain a group representation rather than
an opposite representation.  It is faithful because all coordinate
functions belong to $V_s$.

Let $\mathcal I$ denote the algebra of linear operators on $V_s$ which
send each $V_d$ into $V_{d-1}$.  Then $\mathcal I^{s+1}=0$, and the
representation lies in $1+\mathcal I$.  The elementary identities for
inverses in a nilpotent algebra give
$[1+\mathcal I^a,1+\mathcal I^b]\subseteq1+\mathcal I^{a+b}$.
Thus every $(s+1)$-fold group commutator is the identity.  Extra
translation coordinates can all be given weight one, which proves the
direct-product assertion as well.

For a fixed shear $T$, its displacement in a weight-$j$ coordinate is
$O_T(E^{j-1})$ on $B_E$.  Its derivative matrix is block triangular by
weight with identity diagonal, so its Jacobian determinant is one.
The image $T(B_E)$ is contained in the box with side bounds
$E^j+C_T E^{j-1}$.  The latter box exceeds $B_E$ in relative volume
$O_T(E^{-1})$.  Since $T$ preserves Lebesgue measure, the symmetric
difference of $T(B_E)$ and $B_E$ has the same bound.  This proves
\eqref{eq:fixed-shear-tv}, uniformly over a fixed finite list by taking
the largest of its constants.  The assertion is deliberately for fixed
shears; no uniform coefficient bound over all shears is needed.

For $a\in H$ fixed, left multiplication in logarithmic coordinates is
\[
 u\longmapsto\log(a\exp u).
\]
BCH makes this another shear: every displacement bracket uses the fixed
vector $\log a$ and hence has input weighted degree strictly below its
output weight.  Equation~\eqref{eq:fixed-shear-tv} therefore implies
that every weak subsequential limit of the projected boxes on
$H/(H\cap\sigma\Lambda\sigma^{-1})$ is invariant under all left
translations by $H$.  This compact homogeneous space has a unique
such probability measure, its Haar probability.  Compactness supplies
subsequential limits, and uniqueness shows that the entire sequence
converges.  Translating on the left by $h$ and on the right by
$\sigma\Lambda$ gives the stated point-coset Haar law.  Applying
\eqref{eq:fixed-shear-tv} once more proves the identical limit after
each fixed shear or word.
\end{proof}

The locality and finiteness qualifications in this Section are useful in
Section~\ref{sec:alignment}.  There the marginal face symmetries will be
proved for every fixed residue modulus before $d_0$ is selected.  The
finite recurrence argument will then use only finitely many fixed words
on each limiting coset.  Lemma~\ref{lem:shear-haar} supplies their common
Haar limit without imposing either a dimension bound on those covers or
measurable dependence of the lifts on the limiting coset.
